\subsection{10.03.2026}

\subsubsection{Ziel}

Ziel des Tages war es, den JOY-PI mit dem Internet zu verbinden,
um anschließend die neuen Funktionen der Taupunktberechnung
testen zu können.

\subsubsection{Durchgeführte Arbeiten}

Zunächst wurde versucht, eine Internetverbindung für den
JOY-PI herzustellen, um auf das GitHub-Repository zugreifen
und die aktuellen Änderungen herunterladen zu können.

Anschließend wurden die neuesten Änderungen aus dem
GitHub-Repository gepullt und getestet.

Weitere Änderungen am Projekt konnten während der
Unterrichtseinheit jedoch nicht vorgenommen werden,
da vom Laptop aus zunächst keine Commits möglich waren.

\subsubsection{Problem 1: Fehlende Internetverbindung}

Aufgrund des verwendeten Klassenraums, welcher nicht
ordnungsgemäß in das Schulnetzwerk eingebunden war,
konnte keine Verbindung zum Schulnetzwerk hergestellt werden.

Dadurch war zunächst kein Zugriff auf das
GitHub-Repository möglich.

\textbf{Lösung}

Es wurde ein Hotspot über ein Mobiltelefon eingerichtet,
um eine Internetverbindung herzustellen und die Änderungen
vom GitHub-Repository herunterladen zu können.

\subsubsection{Problem 2: Keine Git-Commits möglich}

Vom Laptop aus war es zunächst nicht möglich,
Änderungen im GitHub-Repository zu committen.

\textbf{Lösung}

Nach dem Unterricht konnte das Problem nach mehreren
Versuchen behoben werden und Commits waren wieder möglich.

Die genaue Ursache des Problems konnte jedoch nicht
eindeutig festgestellt werden.

\subsubsection{Nächste Schritte}

In der nächsten Woche soll der zweite DHT11-Sensor
angeschlossen und in die Taupunktberechnung integriert
werden.

Außerdem sollen die Berechnungsergebnisse in die
Hauptlogik übernommen und anschließend auf dem
LCD-Display ausgegeben werden.